% =====================================================================
% Note on: Treasury Lock Model (anonymous, 4 pages)
% =====================================================================
\documentclass[11pt,a4paper]{article}

\newcommand{\papertag}{anon\_treasury\_lock\_model}

% =====================================================================
% Shared preamble for paper notes
% Include with:  % =====================================================================
% Shared preamble for paper notes
% Include with:  % =====================================================================
% Shared preamble for paper notes
% Include with:  \input{../_common/preamble}
% =====================================================================

% --- Core packages ---
\usepackage[utf8]{inputenc}
\usepackage[T1]{fontenc}
\usepackage{lmodern}
\usepackage[margin=2.5cm]{geometry}
\usepackage{amsmath,amssymb,amsthm,mathtools}
\usepackage{bm}
\usepackage{booktabs}
\usepackage{array}
\usepackage{enumitem}
\usepackage{hyperref}
\usepackage{xcolor}
\usepackage{listings}
\usepackage{fancyhdr}
\usepackage{titlesec}

% --- Hyperref styling ---
\hypersetup{
  colorlinks=true,
  linkcolor=blue!50!black,
  urlcolor=blue!50!black,
  citecolor=blue!50!black
}

% --- Theorem-like environments ---
\theoremstyle{definition}
\newtheorem{defn}{Definition}
\newtheorem{remark}{Remark}
\newtheorem{example}{Example}

\theoremstyle{plain}
\newtheorem{prop}{Proposition}
\newtheorem{lemma}{Lemma}

% --- Coloured boxes for the structured sections ---
% Validation block, commentary, TODOs use coloured frames so the eye
% can find them when flipping through the note.
\usepackage[most]{tcolorbox}

\newtcolorbox{commentary}{
  colback=yellow!5,
  colframe=yellow!50!black,
  title=Commentary,
  fonttitle=\bfseries,
  breakable
}

\newtcolorbox{validation}{
  colback=green!3,
  colframe=green!50!black,
  title=Validation,
  fonttitle=\bfseries,
  breakable
}

\newtcolorbox{todobox}{
  colback=red!3,
  colframe=red!50!black,
  title=TODO / Open questions,
  fonttitle=\bfseries,
  breakable
}

% --- Code listing style for Python snippets in validation blocks ---
\lstdefinestyle{py}{
  language=Python,
  basicstyle=\ttfamily\footnotesize,
  keywordstyle=\color{blue!60!black},
  commentstyle=\color{gray},
  stringstyle=\color{red!60!black},
  showstringspaces=false,
  breaklines=true,
  frame=single,
  framesep=4pt,
  rulecolor=\color{gray!30}
}

% --- Page header: shows paper tag so a printed note is identifiable ---
\pagestyle{fancy}
\fancyhf{}
\fancyhead[L]{\small\textit{\papertag}}
\fancyhead[R]{\small\thepage}
\renewcommand{\headrulewidth}{0.4pt}

% --- Section spacing: tighter than article default ---
\titlespacing*{\section}{0pt}{1.5ex plus 0.5ex minus 0.2ex}{1ex plus 0.2ex}
\titlespacing*{\subsection}{0pt}{1.2ex plus 0.3ex minus 0.2ex}{0.6ex plus 0.2ex}

% =====================================================================

% --- Core packages ---
\usepackage[utf8]{inputenc}
\usepackage[T1]{fontenc}
\usepackage{lmodern}
\usepackage[margin=2.5cm]{geometry}
\usepackage{amsmath,amssymb,amsthm,mathtools}
\usepackage{bm}
\usepackage{booktabs}
\usepackage{array}
\usepackage{enumitem}
\usepackage{hyperref}
\usepackage{xcolor}
\usepackage{listings}
\usepackage{fancyhdr}
\usepackage{titlesec}

% --- Hyperref styling ---
\hypersetup{
  colorlinks=true,
  linkcolor=blue!50!black,
  urlcolor=blue!50!black,
  citecolor=blue!50!black
}

% --- Theorem-like environments ---
\theoremstyle{definition}
\newtheorem{defn}{Definition}
\newtheorem{remark}{Remark}
\newtheorem{example}{Example}

\theoremstyle{plain}
\newtheorem{prop}{Proposition}
\newtheorem{lemma}{Lemma}

% --- Coloured boxes for the structured sections ---
% Validation block, commentary, TODOs use coloured frames so the eye
% can find them when flipping through the note.
\usepackage[most]{tcolorbox}

\newtcolorbox{commentary}{
  colback=yellow!5,
  colframe=yellow!50!black,
  title=Commentary,
  fonttitle=\bfseries,
  breakable
}

\newtcolorbox{validation}{
  colback=green!3,
  colframe=green!50!black,
  title=Validation,
  fonttitle=\bfseries,
  breakable
}

\newtcolorbox{todobox}{
  colback=red!3,
  colframe=red!50!black,
  title=TODO / Open questions,
  fonttitle=\bfseries,
  breakable
}

% --- Code listing style for Python snippets in validation blocks ---
\lstdefinestyle{py}{
  language=Python,
  basicstyle=\ttfamily\footnotesize,
  keywordstyle=\color{blue!60!black},
  commentstyle=\color{gray},
  stringstyle=\color{red!60!black},
  showstringspaces=false,
  breaklines=true,
  frame=single,
  framesep=4pt,
  rulecolor=\color{gray!30}
}

% --- Page header: shows paper tag so a printed note is identifiable ---
\pagestyle{fancy}
\fancyhf{}
\fancyhead[L]{\small\textit{\papertag}}
\fancyhead[R]{\small\thepage}
\renewcommand{\headrulewidth}{0.4pt}

% --- Section spacing: tighter than article default ---
\titlespacing*{\section}{0pt}{1.5ex plus 0.5ex minus 0.2ex}{1ex plus 0.2ex}
\titlespacing*{\subsection}{0pt}{1.2ex plus 0.3ex minus 0.2ex}{0.6ex plus 0.2ex}

% =====================================================================

% --- Core packages ---
\usepackage[utf8]{inputenc}
\usepackage[T1]{fontenc}
\usepackage{lmodern}
\usepackage[margin=2.5cm]{geometry}
\usepackage{amsmath,amssymb,amsthm,mathtools}
\usepackage{bm}
\usepackage{booktabs}
\usepackage{array}
\usepackage{enumitem}
\usepackage{hyperref}
\usepackage{xcolor}
\usepackage{listings}
\usepackage{fancyhdr}
\usepackage{titlesec}

% --- Hyperref styling ---
\hypersetup{
  colorlinks=true,
  linkcolor=blue!50!black,
  urlcolor=blue!50!black,
  citecolor=blue!50!black
}

% --- Theorem-like environments ---
\theoremstyle{definition}
\newtheorem{defn}{Definition}
\newtheorem{remark}{Remark}
\newtheorem{example}{Example}

\theoremstyle{plain}
\newtheorem{prop}{Proposition}
\newtheorem{lemma}{Lemma}

% --- Coloured boxes for the structured sections ---
% Validation block, commentary, TODOs use coloured frames so the eye
% can find them when flipping through the note.
\usepackage[most]{tcolorbox}

\newtcolorbox{commentary}{
  colback=yellow!5,
  colframe=yellow!50!black,
  title=Commentary,
  fonttitle=\bfseries,
  breakable
}

\newtcolorbox{validation}{
  colback=green!3,
  colframe=green!50!black,
  title=Validation,
  fonttitle=\bfseries,
  breakable
}

\newtcolorbox{todobox}{
  colback=red!3,
  colframe=red!50!black,
  title=TODO / Open questions,
  fonttitle=\bfseries,
  breakable
}

% --- Code listing style for Python snippets in validation blocks ---
\lstdefinestyle{py}{
  language=Python,
  basicstyle=\ttfamily\footnotesize,
  keywordstyle=\color{blue!60!black},
  commentstyle=\color{gray},
  stringstyle=\color{red!60!black},
  showstringspaces=false,
  breaklines=true,
  frame=single,
  framesep=4pt,
  rulecolor=\color{gray!30}
}

% --- Page header: shows paper tag so a printed note is identifiable ---
\pagestyle{fancy}
\fancyhf{}
\fancyhead[L]{\small\textit{\papertag}}
\fancyhead[R]{\small\thepage}
\renewcommand{\headrulewidth}{0.4pt}

% --- Section spacing: tighter than article default ---
\titlespacing*{\section}{0pt}{1.5ex plus 0.5ex minus 0.2ex}{1ex plus 0.2ex}
\titlespacing*{\subsection}{0pt}{1.2ex plus 0.3ex minus 0.2ex}{0.6ex plus 0.2ex}

% =====================================================================
% House notation: shared symbols used across all paper notes.
% Each note imports this and may shadow individual macros if a paper
% uses a clashing convention — but the *default* meaning is fixed here.
%
% Include with:  % =====================================================================
% House notation: shared symbols used across all paper notes.
% Each note imports this and may shadow individual macros if a paper
% uses a clashing convention — but the *default* meaning is fixed here.
%
% Include with:  % =====================================================================
% House notation: shared symbols used across all paper notes.
% Each note imports this and may shadow individual macros if a paper
% uses a clashing convention — but the *default* meaning is fixed here.
%
% Include with:  \input{../_common/notation}
% =====================================================================

% --- Measures and expectations ---
\newcommand{\Q}{\mathbb{Q}}              % risk-neutral measure
\newcommand{\Qt}{\mathbb{Q}^{T}}         % T-forward measure
\newcommand{\E}{\mathbb{E}}              % expectation
\newcommand{\Et}[1]{\E^{#1}}             % expectation under measure #1
\newcommand{\Ftime}[1]{\mathcal{F}_{#1}} % filtration at #1

% --- Discount and forward objects ---
\newcommand{\Disc}[2]{D_{#1,#2}}         % discount factor D_{t,T}
\newcommand{\Bond}[2]{P_{#1}(#2)}        % bond price P_t(y)
\newcommand{\Ann}[1]{A_{#1}}             % annuity / level
\newcommand{\Numer}[1]{N_{#1}}           % numeraire
\newcommand{\Fwd}[2]{F_{#1,#2}}          % forward price F_{t,T}
\newcommand{\Fwdpx}[2]{\mathrm{Fwd}_{#1}(#2)} % verbose forward-price F_t(T)

% --- Rates ---
\newcommand{\IRR}{\mathrm{IRR}}          % internal rate of return
\newcommand{\Yld}{y}                     % bond yield variable
\newcommand{\Strike}{K}                  % strike (price-space)
\newcommand{\Lock}{L}                    % locked yield / rate
\newcommand{\Repo}{r^{\mathrm{repo}}}    % repo rate
\newcommand{\Fund}{r^{\mathrm{fun}}}     % unsecured funding rate
\newcommand{\OIS}{r^{\mathrm{ois}}}      % OIS rate
\newcommand{\Coll}{r^{\mathrm{c}}}       % collateral rate
\newcommand{\Rcmt}{R^{\mathrm{cmt}}}     % constant-maturity-treasury rate
\newcommand{\Rswp}{R^{\mathrm{swp}}}     % swap rate (used for risky variant)
\newcommand{\Rec}{\mathrm{Rec}}          % recovery rate
\newcommand{\NPV}{\mathrm{NPV}}          % present-value functional
\newcommand{\PVone}{\mathrm{PV01}}       % risky annuity
\newcommand{\Loss}{\mathrm{Loss}}        % default-leg integral
\newcommand{\Sasw}{S^{\mathrm{ASW}}}     % asset-swap spread
\newcommand{\Scds}{S^{\mathrm{CDS}}}     % CDS spread
\newcommand{\Basis}{\mathrm{Basis}}      % CDS-bond basis

% --- Sensitivities ---
\newcommand{\Diff}[2]{D_{#1}^{#2}}       % Euler's notation: D_x^k[f]
\newcommand{\Riskf}{\mathrm{RiskFactor}} % bond risk factor (= -DV01*1e4)
\newcommand{\DVone}{\mathrm{DV01}}

% --- Indicator / misc ---
\newcommand{\Ind}{\mathbf{1}}
\newcommand{\given}{\,|\,}
\newcommand{\dd}{\,\mathrm{d}}

% --- Greeks-style ---
\newcommand{\Gam}{\Gamma}
\newcommand{\Vega}{\mathcal{V}}

% =====================================================================

% --- Measures and expectations ---
\newcommand{\Q}{\mathbb{Q}}              % risk-neutral measure
\newcommand{\Qt}{\mathbb{Q}^{T}}         % T-forward measure
\newcommand{\E}{\mathbb{E}}              % expectation
\newcommand{\Et}[1]{\E^{#1}}             % expectation under measure #1
\newcommand{\Ftime}[1]{\mathcal{F}_{#1}} % filtration at #1

% --- Discount and forward objects ---
\newcommand{\Disc}[2]{D_{#1,#2}}         % discount factor D_{t,T}
\newcommand{\Bond}[2]{P_{#1}(#2)}        % bond price P_t(y)
\newcommand{\Ann}[1]{A_{#1}}             % annuity / level
\newcommand{\Numer}[1]{N_{#1}}           % numeraire
\newcommand{\Fwd}[2]{F_{#1,#2}}          % forward price F_{t,T}
\newcommand{\Fwdpx}[2]{\mathrm{Fwd}_{#1}(#2)} % verbose forward-price F_t(T)

% --- Rates ---
\newcommand{\IRR}{\mathrm{IRR}}          % internal rate of return
\newcommand{\Yld}{y}                     % bond yield variable
\newcommand{\Strike}{K}                  % strike (price-space)
\newcommand{\Lock}{L}                    % locked yield / rate
\newcommand{\Repo}{r^{\mathrm{repo}}}    % repo rate
\newcommand{\Fund}{r^{\mathrm{fun}}}     % unsecured funding rate
\newcommand{\OIS}{r^{\mathrm{ois}}}      % OIS rate
\newcommand{\Coll}{r^{\mathrm{c}}}       % collateral rate
\newcommand{\Rcmt}{R^{\mathrm{cmt}}}     % constant-maturity-treasury rate
\newcommand{\Rswp}{R^{\mathrm{swp}}}     % swap rate (used for risky variant)
\newcommand{\Rec}{\mathrm{Rec}}          % recovery rate
\newcommand{\NPV}{\mathrm{NPV}}          % present-value functional
\newcommand{\PVone}{\mathrm{PV01}}       % risky annuity
\newcommand{\Loss}{\mathrm{Loss}}        % default-leg integral
\newcommand{\Sasw}{S^{\mathrm{ASW}}}     % asset-swap spread
\newcommand{\Scds}{S^{\mathrm{CDS}}}     % CDS spread
\newcommand{\Basis}{\mathrm{Basis}}      % CDS-bond basis

% --- Sensitivities ---
\newcommand{\Diff}[2]{D_{#1}^{#2}}       % Euler's notation: D_x^k[f]
\newcommand{\Riskf}{\mathrm{RiskFactor}} % bond risk factor (= -DV01*1e4)
\newcommand{\DVone}{\mathrm{DV01}}

% --- Indicator / misc ---
\newcommand{\Ind}{\mathbf{1}}
\newcommand{\given}{\,|\,}
\newcommand{\dd}{\,\mathrm{d}}

% --- Greeks-style ---
\newcommand{\Gam}{\Gamma}
\newcommand{\Vega}{\mathcal{V}}

% =====================================================================

% --- Measures and expectations ---
\newcommand{\Q}{\mathbb{Q}}              % risk-neutral measure
\newcommand{\Qt}{\mathbb{Q}^{T}}         % T-forward measure
\newcommand{\E}{\mathbb{E}}              % expectation
\newcommand{\Et}[1]{\E^{#1}}             % expectation under measure #1
\newcommand{\Ftime}[1]{\mathcal{F}_{#1}} % filtration at #1

% --- Discount and forward objects ---
\newcommand{\Disc}[2]{D_{#1,#2}}         % discount factor D_{t,T}
\newcommand{\Bond}[2]{P_{#1}(#2)}        % bond price P_t(y)
\newcommand{\Ann}[1]{A_{#1}}             % annuity / level
\newcommand{\Numer}[1]{N_{#1}}           % numeraire
\newcommand{\Fwd}[2]{F_{#1,#2}}          % forward price F_{t,T}
\newcommand{\Fwdpx}[2]{\mathrm{Fwd}_{#1}(#2)} % verbose forward-price F_t(T)

% --- Rates ---
\newcommand{\IRR}{\mathrm{IRR}}          % internal rate of return
\newcommand{\Yld}{y}                     % bond yield variable
\newcommand{\Strike}{K}                  % strike (price-space)
\newcommand{\Lock}{L}                    % locked yield / rate
\newcommand{\Repo}{r^{\mathrm{repo}}}    % repo rate
\newcommand{\Fund}{r^{\mathrm{fun}}}     % unsecured funding rate
\newcommand{\OIS}{r^{\mathrm{ois}}}      % OIS rate
\newcommand{\Coll}{r^{\mathrm{c}}}       % collateral rate
\newcommand{\Rcmt}{R^{\mathrm{cmt}}}     % constant-maturity-treasury rate
\newcommand{\Rswp}{R^{\mathrm{swp}}}     % swap rate (used for risky variant)
\newcommand{\Rec}{\mathrm{Rec}}          % recovery rate
\newcommand{\NPV}{\mathrm{NPV}}          % present-value functional
\newcommand{\PVone}{\mathrm{PV01}}       % risky annuity
\newcommand{\Loss}{\mathrm{Loss}}        % default-leg integral
\newcommand{\Sasw}{S^{\mathrm{ASW}}}     % asset-swap spread
\newcommand{\Scds}{S^{\mathrm{CDS}}}     % CDS spread
\newcommand{\Basis}{\mathrm{Basis}}      % CDS-bond basis

% --- Sensitivities ---
\newcommand{\Diff}[2]{D_{#1}^{#2}}       % Euler's notation: D_x^k[f]
\newcommand{\Riskf}{\mathrm{RiskFactor}} % bond risk factor (= -DV01*1e4)
\newcommand{\DVone}{\mathrm{DV01}}

% --- Indicator / misc ---
\newcommand{\Ind}{\mathbf{1}}
\newcommand{\given}{\,|\,}
\newcommand{\dd}{\,\mathrm{d}}

% --- Greeks-style ---
\newcommand{\Gam}{\Gamma}
\newcommand{\Vega}{\mathcal{V}}


\title{Note on: \emph{Treasury Lock Model}\\
       \large Anonymous model documentation}
\author{Reading notes}
\date{\today}

\begin{document}
\maketitle

% =====================================================================
\section*{Metadata}
\begin{tabular}{@{}p{3.5cm}p{11cm}@{}}
\toprule
Title           & Treasury Lock Model \\
Author(s)       & Not stated \\
Date            & Not stated \\
Source          & Document references \url{finpricing.com}; no SSRN/DOI \\
Local file      & \texttt{papers/anon\_treasury\_lock\_model/31670362145195.pdf} \\
Tags            & \texttt{treasury}, \texttt{t-lock}, \texttt{forward},
                  \texttt{repo}, \texttt{pv01}, \texttt{model-doc} \\
Date read       & \today \\
\bottomrule
\end{tabular}

% =====================================================================
\section*{One-paragraph summary}
A four-page model-documentation note describing the implementation of a
T-Lock pricer. It states the bond forward price as a clean-price formula
discounted with the repo curve, then writes the T-Lock NPV in three
equivalent forms depending on what is locked: the clean price, the
yield (via PV01), or the dirty price. The yield-lock form uses a
discrete two-sided PV01 around the forward yield rather than a
closed-form derivative. The document closes with a caveat that PV01
and NDELTA are underlier-related and only approximate the T-Lock's own
sensitivities.

% =====================================================================
\section{Setup and notation}

The document mixes price-quote conventions (clean/dirty/forward) with
yield-locked conventions. It uses one repo discount factor for the
forward-price construction and a separate ``bond zero curve''
discount factor for the T-Lock NPV. This dual-curve structure is not
typical of more academic treatments (e.g.\ Pucci 2019, where a single
risk-free DF $\Disc{t}{t_e}$ is used and the repo enters only in the
forward drift).

\begin{tabular}{@{}p{3cm}p{3.5cm}p{8.5cm}@{}}
\toprule
Document         & House                  & Meaning \\
\midrule
$P_0$            & $P^{mkt,\text{clean}}_t$ & Clean price of the underlying bond at spot \\
$A_0, A_{del}$   & ---                     & Accrued interest at spot, at delivery \\
$T_0, T_{del}$   & $t$, $t_e$              & Spot settlement, forward settlement (delivery) dates \\
$N$              & ---                     & Number of coupons between $T_0$ and $T_{del}$ \\
$C_i, T_i$       & $c_i, t_i$              & $i$-th coupon and its payment date \\
$B_f$            & $\Fwdpx{t}{t_e}$ (clean) & Bond forward (clean) price at delivery \\
$D_{repo}(\cdot,\cdot)$ & ---              & Discount factor from the repo zero curve \\
$D_{\text{T-Lock}}(\cdot,\cdot)$ & ---     & Discount factor from the bond zero curve (used to discount T-Lock NPV) \\
$M$              & $N$ (notional)          & Trade amount \\
$y_f$            & $\IRR_{\text{fwd}}$     & Forward bond yield at delivery \\
$y_{\text{T-Lock}}$ & $\Lock$              & Locked yield \\
$P_{\text{T-Lock}}$ & $\Lock^{(P)}$        & Locked clean price \\
$PV01(y_f)$      & ---                     & Two-sided finite-difference PV01 at $y_f$ \\
$B(y)$           & $\Bond{t}{y}$           & Bond PV at yield $y$ \\
\bottomrule
\end{tabular}

% =====================================================================
\section{Key equations}

\paragraph{Bond forward (clean) price.}
\begin{equation}\label{eq:fwd-clean}
  B_f \;=\; \frac{P_0 + A_0}{D_{repo}(T_0, T_{del})}
          \;-\; \sum_{i=1}^{N} \frac{C_i}{D_{repo}(T_i, T_{del})}
          \;-\; A_{del}.
\end{equation}
\textit{Says:} accrete the dirty spot price $(P_0 + A_0)$ to delivery
at the repo rate, subtract coupon flows that occur in $[T_0, T_{del}]$
each accreted to delivery, then strip the delivery-date accrued
interest to land on a \emph{clean} forward. This is the same content
as Pucci's eq.~(49) in continuous form, here in discrete-DF form, and
expressed clean instead of dirty.

\paragraph{T-Lock NPV: clean-price lock.}
\begin{equation}\label{eq:npv-clean}
  \mathrm{NPV}_{\text{T-Lock}} \;=\; N \cdot \bigl(P_{\text{T-Lock}} - B_f\bigr) \cdot D_{\text{T-Lock}}(t, T_{del}).
\end{equation}
\textit{Says:} a forward contract on the clean price. The reference
forward $B_f$ is computed from \eqref{eq:fwd-clean} using the repo
curve; the resulting payoff is then discounted to $t$ using a separate
``bond zero'' discount curve.

\paragraph{T-Lock NPV: yield lock (via PV01).}
\begin{equation}\label{eq:npv-yield}
  \mathrm{NPV}_{\text{T-Lock}} \;=\; M \cdot (y_f - y_{\text{T-Lock}}) \cdot PV01(y_f) \cdot D_{\text{T-Lock}}(t, T_{del}),
\end{equation}
with the PV01 defined as a two-sided finite difference
\begin{equation}\label{eq:pv01-fd}
  PV01(y_f) \;=\; B(y_f + 0.5\,\text{bp}) - B(y_f - 0.5\,\text{bp}).
\end{equation}
\textit{Says:} the linear payoff $(y_f - L) \cdot \text{PV01}$ is the
discrete analogue of Pucci's $\Riskf_{t_e} \cdot (\IRR_{t_e} - \Lock)$.
The sign convention here makes $PV01(y_f)$ \emph{positive} (because a
positive yield shift lowers price, but the formula computes
$B(y_f+) - B(y_f-)$ with $B$ decreasing in $y$, giving a negative
quantity --- see commentary).

\paragraph{T-Lock NPV: dirty-price lock.}
The document states that the dirty case is obtained from the
clean-price case via $P^{dirty} = P^{clean} + A$, i.e.\ it computes the
clean-price lock and adjusts by accrued interest. No new formula.

% =====================================================================
\section{Derivations \& commentary}

\begin{commentary}
The PV01 sign convention in \eqref{eq:pv01-fd} as written gives a
\emph{negative} number (because $B$ is decreasing in $y$). If one
takes the document's formula literally then a long T-Lock with $y_f >
y_{\text{T-Lock}} = \Lock$ produces a negative NPV, which contradicts
the intent (long T-Lock pays when yields rise). The conventional
practice is either (a) define PV01 as the absolute price change per
$+1$~bp, or (b) flip the sign of $(y_f - \Lock)$. The document does
not say which it adopts. House interpretation: take
$\widetilde{PV01}(y_f) = -[B(y_f + 0.5\text{bp}) - B(y_f - 0.5\text{bp})] > 0$,
which matches Pucci's $\Riskf_{t_e} = -\partial \Bond{}{y}/\partial y > 0$
in the limit of small bump width.
\end{commentary}

\begin{commentary}
The document uses two distinct discount curves: a repo curve for the
forward bond price, and a ``bond zero curve'' for the T-Lock NPV.
The choice of the latter is unusual; Pucci (2019) and standard
post-2008 practice would use a collateral/OIS rate for derivative
discounting, not a bond curve. This is the kind of detail that should
be confirmed against the firm's actual CSA before using the formula.
\end{commentary}

\begin{commentary}
The finite-difference PV01 in \eqref{eq:pv01-fd} differs from a
closed-form $-\partial \Bond{}{y}/\partial y$ by terms of order
$\text{bp}^2$. For practical purposes ($1$~bp is small) the difference
is negligible, but a property test should compare the two and check
that the finite-difference converges to the analytic derivative as
the bump shrinks.
\end{commentary}

% =====================================================================
\section{Validation}

\begin{validation}
This is a \emph{specification} of mathematical objects the paper
defines and the numerical values they should take. It says nothing
about how those objects are implemented or invoked.

\textbf{Quantities the paper defines:}
\begin{itemize}[noitemsep]
  \item $B_f$ in \eqref{eq:fwd-clean}, the clean forward bond price
        constructed from the repo-curve discount factors.
  \item $\mathrm{NPV}_{\text{T-Lock}}$ in three equivalent
        formulations: clean-price lock \eqref{eq:npv-clean},
        yield lock \eqref{eq:npv-yield}, dirty-price lock (derived).
  \item $PV01(y_f)$ as a two-sided finite difference
        \eqref{eq:pv01-fd}.
\end{itemize}

\textbf{Canonical numerical case:}
The document provides no worked example. For validation purposes,
re-use the Pucci 2019 canonical case (US Treasury 3.125\% 25-Nov-2028,
ISIN US9128285M81, 24-Jan-2019, $t_e - t = 3$ months,
$P^{mkt,dirty} = 104.1055$, $\IRR = 2.717\%$, $\Repo = 2.46\%$); the
two documents should agree to first order.

\textbf{Expected numerical values:} not available in this document.
Cross-check is against the Pucci numbers ($B_f$ dirty $= 104.74$;
forward IRR $= 2.53\%$) within $\pm 0.01$ and $\pm 1$~bp.

\textbf{Equation $\leftrightarrow$ quantity map:}

\smallskip
\begin{tabular}{@{}p{2.8cm}p{3.8cm}p{6.2cm}@{}}
\toprule
Paper eq.                & Symbol                                & Quantity \\
\midrule
\eqref{eq:fwd-clean}     & $B_f$                                 & clean forward bond price via repo discount factors \\
\eqref{eq:npv-clean}     & $\mathrm{NPV}_{\text{T-Lock}}$        & T-Lock NPV when the clean price is locked \\
\eqref{eq:npv-yield}     & $\mathrm{NPV}_{\text{T-Lock}}$        & T-Lock NPV when the yield is locked, via PV01 \\
\eqref{eq:pv01-fd}       & $PV01(y_f)$                           & two-sided finite-difference PV01 at the forward yield \\
\bottomrule
\end{tabular}
\smallskip

\textbf{Invariants and identities to verify:}
\begin{enumerate}[noitemsep]
  \item \emph{Equivalence of clean and dirty locks}: the dirty-price
        formulation, after subtracting $A_{del}$ from the locked
        price, equals the clean-price formulation
        \eqref{eq:npv-clean}.
  \item \emph{First-order agreement with the analytic derivative}: as
        the FD bump width $\to 0$, $PV01(y_f) / \text{bp}$ converges
        to $-\partial \Bond{}{y} / \partial y |_{y=y_f}$.
  \item \emph{First-order agreement of yield and clean-price locks}:
        for $y_{\text{T-Lock}}$ near $y_f$, the clean-price and yield
        formulations agree to $O((y_f - y_{\text{T-Lock}})^2)$ when
        the locked clean price is $\Bond{t_e}{y_{\text{T-Lock}}}$.
  \item \emph{Repo-curve consistency}: $B_f$ as constructed in
        \eqref{eq:fwd-clean} satisfies the no-arbitrage condition that
        a repo strategy returns zero P\&L when valued at $T_{del}$.
\end{enumerate}
\end{validation}

% =====================================================================
\section{Glossary of symbols (paper's notation)}
\begin{tabular}{@{}p{3cm}p{12cm}@{}}
\toprule
$P_0$                                 & Clean price of the underlying bond at spot \\
$A_0$                                 & Accrued interest on a unit underlying bond at $T_0$ \\
$T_0$                                 & Spot settlement date \\
$T_{del}$                             & Delivery date (forward settlement date) \\
$A_{del}$                             & Accrued interest on a unit underlying bond at $T_{del}$ \\
$N$ (eq.~\ref{eq:fwd-clean})          & Number of coupons paid in $[T_0, T_{del}]$ \\
$N$ (eq.~\ref{eq:npv-clean})          & Quantity of the underlying bond (notional) \\
$C_i$                                 & $i$-th coupon of the unit underlying bond \\
$T_i$                                 & $i$-th coupon payment date \\
$D_{repo}(t_1, t_2)$                  & Discount factor from $t_2$ back to $t_1$, repo zero curve \\
$D_{\text{T-Lock}}(t_1, t_2)$         & Discount factor from $t_2$ back to $t_1$, bond zero yield curve \\
$B_f$                                 & Bond forward clean price at delivery \\
$P_{\text{T-Lock}}$                   & Locked clean price for the underlying bond \\
$M$                                   & Trade amount of the underlying bond \\
$y_f$                                 & Forward bond yield (as if bond settles at delivery) \\
$y_{\text{T-Lock}}$                   & Locked bond yield \\
$PV01(y_f)$                           & $B(y_f + 0.5\,\text{bp}) - B(y_f - 0.5\,\text{bp})$ \\
$B(y)$                                & Present value of the underlying bond given yield $y$ \\
$t$                                   & Valuation date \\
\bottomrule
\end{tabular}

% =====================================================================
\begin{todobox}
\begin{itemize}[noitemsep]
  \item Confirm the PV01 sign convention used in eq.~\ref{eq:npv-yield}
        --- as written, the signs of $(y_f - \Lock)$ and the FD-PV01
        do not produce the conventional ``long T-Lock pays on yield
        rise'' behaviour.
  \item The document uses a separate ``bond zero curve'' for
        discounting the T-Lock NPV. Confirm against firm CSA whether
        this is appropriate; modern practice would normally use
        OIS/collateral.
  \item No worked numerical example in the document. Cross-validate
        \eqref{eq:fwd-clean} against the Pucci 24-Jan-2019 case.
\end{itemize}
\end{todobox}

\end{document}
